\section{Introduction}\label{sec:intro}

At fixed information, a stronger incumbent reduces a prospective challenger's expected return from entry into a
takeover contest. The challenger must pay preparation costs before it can submit an executable offer, and a
stronger rival lowers what that preparation earns \citep{Fishman1988,HirshleiferPng1989}. I study a force in
the opposite direction. Competition can make the target's stock more informative about the challenger, and when a
potential buyer observes that price before it commits to acquisition preparation, the information effect can
outweigh the loss of acquisition profits. A stronger incumbent then attracts another buyer and raises the
probability that a challenger from the high class acquires the target.

The setting is a listed target with a public sale opportunity. One buyer is already prepared, and another has
not committed to the diligence that an executable proposal requires; Section~\ref{sec:interval} describes the
decision interval in which that buyer decides while the target's stock trades. Incumbent strength is the public
distribution of the prepared bidder's value; a stronger incumbent shifts it upward, and the challenger decides
before any offer is announced.

The challenger and target shareholders hold two claims on the same acquisition surplus: the challenger keeps
what remains if it prepares, bids, and wins, and target shareholders hold the traded claim, what the cash
second-price auction pays them. The challenger's value comes from one of two quality classes, high and low, and
the classes fix the expected value gap between them. A stronger incumbent outbids the challenger more often and
then pays the challenger's bid, so more of the gap reaches target proceeds; a challenger that wins pays more, so
less of the gap remains as acquisition profit at every belief. Target shares therefore become more sensitive to
information about the challenger at the same time as the challenger's own return to preparation falls.

An outside specialist investor with no net position in the target knows the challenger's class and trades on that
sensitivity. The trading stage is one order slot, with the alternative-arrival structure of
\citet{DowGorton1997}: with some probability the investor fills the slot, and otherwise a liquidity trader fills
it with a sale, a purchase, or no order. No one observes who placed the order, and a liquidity order can match
the investor's direction, so no order reveals the class. Under every investor strategy the posterior after every
order stays inside fixed bounds, and the investor keeps a residual advantage after every order; a larger gap in
target proceeds can make that advantage worth the cost of trade. Target shareholders gain when the price recruits
entry, because a holder through the sale receives sale proceeds, not the interim price.

The main result compares three economies on a nonempty open set of parameters: one with a weak incumbent, one
with a strong incumbent, and one with a very strong incumbent against which recruitment stops. In each economy the
equilibrium outcome is unique. With the weak incumbent, the investor's residual advantage cannot
cover the trading cost, the price reveals nothing, and only a zero-cost challenger prepares. With the strong
incumbent, the investor buys after good news and sells after bad news, and a purchase brings in a high-cost
challenger, so entry rises even though acquisition profit falls at every fixed belief, and under a condition on
the high class's chance to win, which the benchmark meets, the probability that a
challenger from the high-value class owns the target rises with it. With the very strong incumbent, the price
stays informative, but even the most favorable price fails to justify preparation at the high cost, so entry
returns to its floor.

The reversal comes from the change in information that competition induces. Market makers anticipate the
preparation that each price induces, so the reversal holds under rational pricing; when the investor's orders
stay fixed at their informative level while the incumbent strengthens, a stronger incumbent can only lower entry
(Section~\ref{sec:benchmark}).

In the benchmark of Section~\ref{sec:benchmark}, entry rises from \reg{entry_r0} at the weak incumbent to
\reg{entry_r1_B} at the strong incumbent, and high-class ownership rises from \reg{own_r0} to \reg{own_r1_B}.
Over the same change, gross acquisition profit at the prior falls from \reg{Bprior_r0} to \reg{Bprior_r1_B},
while the target-payoff spread rises from \reg{Delta_r0} to \reg{Delta_r1}. In the
benchmark, high-cost preparation stops above the entry ceiling, the strength \(r_C\approx\reg{rC_B}\). At the
benchmark auction parameters, the hypotheses of the main result hold exactly when
\(\reg{u_lo_frac}<u<\reg{u_hi_frac}\), where \(u\) is the informed presence mass relative to the liquidity mass
in one direction. Equivalently, the top posterior at the prior \(1/2\) lies between the preparation thresholds
of the strong and the very strong incumbent, \(\reg{M_lo_frac}\) and \(\reg{M_hi_frac}\).
Section~\ref{sec:benchmark} reports this region and the range of priors in a table.

Informative and uninformative equilibria coexist on an interval of strengths
(Section~\ref{sec:coexistence}), and Appendix~\ref{app:coexistence} lists the complete equilibrium set along
the benchmark path. In
the benchmark that interval runs from \(r_A\approx\reg{rA}\) to \(r_N\approx\reg{rN}\). Along the informative
branch, when the class is high the investor first mixes between a purchase and abstention, and entry rises with
strength; from \(r_F\approx\reg{rF_B}\), the investor buys for sure in that class, and entry stays constant on
the rest of the interval. The same start of coexistence holds for a window of priors around the benchmark.

When the investor pays for research that can make it informed, research and entry rise together at intermediate
strengths (Section~\ref{sec:research}). The investor chooses the probability of success up to a cap, at a
quadratic cost. The unique outcome has no research with the weak incumbent and research at the cap with the
strong incumbent. At the strengths \reg{str_r155}, \reg{str_r160} and \reg{str_r165}, research rises from
\reg{res_lambda_r155} through \reg{res_lambda_r160} to \reg{res_lambda_r165}, and entry from
\reg{res_entry_r155} through \reg{res_entry_r160} to \reg{res_entry_r165}. Along this branch the comovement
holds for every liquidity probability of the model, and the branch coexists with an equilibrium without research.

The rise in entry survives other laws of preparation costs, with or without atoms
(Section~\ref{sec:costs}). For every prior and every pair of classes that satisfy the class ranking and differ
above the reserve, the reversal holds at some incumbent strengths and costs (Section~\ref{sec:classes}); with
more than two classes, it holds under a separation condition. It also survives private signals that the challenger
and the investor hold about different facts (Section~\ref{sec:signals}); in an example with its own costs and
strengths, the challenger's signal is more accurate than the investor's, and the rise in that example needs both
signals.

In the strong-incumbent economy, access to prices raises expected target proceeds and acquisition surplus net
of preparation costs (Section~\ref{sec:access}). Sale terms shape both what a buyer pays when it wins and what
the stock price can reveal before anyone prepares. An alternative bargaining institution shows how the
payment rule matters: the effect of competition on the target-payoff spread depends on the weight the rule
places on the runner-up's value (Section~\ref{sec:payment}). A reserve that excludes the low class can make the
price informative against both incumbents, and if informed trade is then full and only a purchase recruits at both strengths,
entry is flat in incumbent strength (Section~\ref{sec:reserves}).

A purchase raises the price by more than a sale lowers it, both measured from the price after no order,
exactly when the purchase recruits the challenger (Section~\ref{sec:empirical}). The gap comes from the entry
component of the spread between purchase and sale prices; in the benchmark, that component is
\reg{ba_share_r1_B} of the spread. Above the entry ceiling, orders no longer predict entry, but they still
predict which bidder wins; given the entry floor and the target-payoff spread, the sale impact and the volume
of orders identify the informed share.

The paper belongs to the literature on the real effects of financial markets \citep{BondEdmansGoldstein2012}.
The closest antecedent is \citet{DowGoldsteinGuembel2017}, where a firm's investment decision feeds back into
the incentive to produce information about its stock; here the sale rule divides acquisition surplus between a
traded claim and a challenger that decides whether to prepare. \citet{EdmansGoldsteinJiang2015} show that corrective
real decisions can discourage trade on bad news under rational pricing; here the real decision is the challenger's
entry itself.

Firms learn from peers' stock prices \citep{FoucaultFresard2014}, and target prices affect takeover activity
\citep{EdmansGoldsteinJiang2012}; in this model, a prospective challenger learns about its own acquisition value
from the target's price before it commits.\footnote{Bidders also learn from announcement returns when they decide
whether to complete a deal \citep{Luo2005}. That decision comes after the entry margin of this paper.}

The auction-entry literature makes the bidder pool endogenous
\citep{Samuelson1985,McAfeeMcMillan1987,LevinSmith1994,GentryStroup2019}. Under selective entry, a bidder's
pre-entry information indexes the distribution of its value, and the payment of the entry cost reveals the value
\citep{RobertsSweeting2013}.\footnote{One limit of this framework is a bidder that knows only the common value
distribution before entry \citep{LevinSmith1994}. The other limit is a bidder that already knows its value
\citep{Samuelson1985}. Selective entry also affects the choice of sale procedure \citep{RobertsSweeting2013}, and
auction formats affect bidders' incentives to acquire information \citep{Persico2000}.} Here the pre-entry
information is the target's stock price, a public price that moves with the orders of an investor who observes
the challenger's quality class. I keep the direct deterrence force of the auction-entry literature and add a
market that operates before preparation, with the informed trader outside the auction and trading a payoff that
differs from the entrant's profit.\footnote{Recent work studies how bidders learn about their own values and
about competitors \citep{PernoudGleyze2026}, post-auction feedback in security-payment design
\citep{LiuBernhardt2022}, and bidder-pool choice with correlated values \citep{CarlinEtAl2026}. In takeovers,
\citet{BettonEtAl2014} study negotiations with stock-market feedback, and \citet{LinMaYangZhu2025} model payment
choice and trade within an initiated deal; I hold cash consideration fixed and study entry.}

The investor in this paper is an outside trader and earns only the trading residual. An informed trader can also hold a stake before a control decision: an acquirer can use
liquidity trade to conceal share purchases before a bid \citep{KyleVila1991}, large holders trade on their
information \citep{Edmans2009}, and arbitrageurs' positions can affect tendering
\citep{CornelliLi2002}.\footnote{Large holders also monitor or intervene as prices feed back to firm decisions
\citetext{\citealp[Section~II.E]{EdmansGoldsteinJiang2015}; \citealp{Maug1998}; \citealp{BackEtAl2018}}.
\citet{GoldsteinGuembel2008} study manipulation through the real decision that responds to the price. The model
abstracts from toeholds and dispersed-shareholder free riding \citep{BulowHuangKlemperer1999,GrossmanHart1980}.}

Section~\ref{sec:model} sets out the model, and Section~\ref{sec:payoffs} derives the returns to information.
Section~\ref{sec:results} states the main results. Sections~\ref{sec:research} to~\ref{sec:empirical} cover research,
robustness, welfare and sale terms, and empirical implications, and Section~\ref{sec:conclusion} concludes.
Appendix~\ref{app:proofs} states the full form of Proposition~\ref{prop:coexistence} and the results of
Sections~\ref{sec:research} to~\ref{sec:empirical} in full, and proves every result.
